\documentclass[11pt]{article}
\usepackage[margin=1in]{geometry}
\usepackage{amsmath,amssymb,amsthm,mathtools}
\usepackage{microtype}
\usepackage[colorlinks=true,linkcolor=blue,urlcolor=blue,citecolor=blue]{hyperref}

\newtheorem{theorem}{Theorem}
\newtheorem{corollary}[theorem]{Corollary}
\newtheorem{proposition}[theorem]{Proposition}
\newcommand{\HS}{\mathcal S_2}
\newcommand{\TC}{\mathcal S_1}
\newcommand{\ran}{\operatorname{Ran}}
\newcommand{\supp}{\operatorname{supp}}

\title{A projection-support formula for Hilbert--Schmidt closure}
\author{Candidate partial answer to Question 4.14 of arXiv:2511.23121}
\date{11 August 2026}

\begin{document}
\maketitle

\noindent\fbox{\parbox{0.94\linewidth}{\textbf{Status: candidate partial
result, likely valid.}  The Hilbert--Schmidt closure projection is exactly
the complement of a join of support projections of trace-class test
vectors.  For $M=\mathcal B(H)$ this completely classifies all
finite-corank projections, including a sharp corank-one dichotomy.  The
formula does not yet eliminate the trace-class test-vector condition for an
arbitrary von Neumann algebra, so it is not claimed as a complete intrinsic
classification in the full generality requested by the source.}}

\section{Source question and notation}

Let $M\subseteq\mathcal B(H)$ be a von Neumann algebra and put
\[
 K=\HS(H),\qquad A=M\,\bar\otimes\,M^{\mathrm{op}}\subseteq\mathcal B(K).
\]
Under the correspondence in Daws's paper, an HS quantum relation is
$V=eK$ for a projection $e\in A$.  Its Hilbert--Schmidt closure is
\[
 V^{\mathrm{hsc}}=(\overline V^{\,w^*})\cap\HS(H),
\]
and corresponds to a projection $e^{\mathrm{hsc}}\in A$ with
$e\leq e^{\mathrm{hsc}}$.  Question~4.14 asks for a description of
$e\mapsto e^{\mathrm{hsc}}$ at the level of
$M\bar\otimes M^{\mathrm{op}}$ \cite[Question~4.14]{Daws}.

Write $D=\TC(H)\subset K$.  For $\xi\in K$, let
\[
 \omega_\xi(a)=\langle a\xi,\xi\rangle\quad(a\in A),
 \qquad s_A(\omega_\xi)\in A
\]
denote the support projection of this normal positive functional.

\section{The support-hull formula}

\begin{theorem}[Projection-support formula]\label{thm:formula}
For every projection $e\in A$,
\begin{equation}\label{eq:main}
 e^{\mathrm{hsc}}
 =1-\bigvee\bigl\{s_A(\omega_\xi):
                 \xi\in\TC(H),\ e\xi=0\bigr\}.
\end{equation}
Equivalently, if
\[
 W_e=\overline{\TC(H)\cap(1-e)K}^{\,\|\cdot\|_2},
\]
then the orthogonal projection $r_e$ onto $W_e$ belongs to $A$ and
\begin{equation}\label{eq:re}
                  e^{\mathrm{hsc}}=1-r_e.
\end{equation}
In particular,
\begin{equation}\label{eq:fixed}
 e=e^{\mathrm{hsc}}quad\Longleftrightarrow\quad
 \overline{\TC(H)\cap\ker e}^{\,\|\cdot\|_2}=\ker e.
\end{equation}
\end{theorem}

\paragraph{Proof intuition.}
The source identifies $V^{\mathrm{hsc}}$ as the Hilbert-space annihilator
of the trace-class vectors normal to $V$.  Thus its orthogonal complement
is the closed trace-class part of $\ker e$.  That subspace is automatically
reducing for the commutant of $A$, so its projection lies in $A$.  Support
projections package the same closed subspace as a join inside the projection
lattice of $A$.

\begin{proof}
Proposition~4.11 of the source gives
\begin{equation}\label{eq:source}
 V^{\mathrm{hsc}}=(\TC(H)\cap V^\perp)^\perp
                 =(\TC(H)\cap(1-e)K)^\perp,
\end{equation}
where both orthogonal complements are taken in $K$.

Let $A'$ be the commutant of the displayed action of $A$ on $K$.
Trace-class operators form a two-sided operator ideal, hence $D$ is
invariant under the left and right bounded multiplications comprising
$A'$.  Since $e\in A$, the space $(1-e)K$ is also $A'$-invariant.
Consequently $W_e$ is invariant under the self-adjoint algebra $A'$ and is
therefore reducing for $A'$.  Its orthogonal projection $r_e$ commutes with
$A'$, so
\[
                         r_e\in(A')'=A.
\]
Equation~\eqref{eq:source} now says that
$\ran(e^{\mathrm{hsc}})=W_e^\perp$, proving
\eqref{eq:re}.

It remains to identify $r_e$ as the join in \eqref{eq:main}.  For any
$\xi\in K$, the projection onto $\overline{A'\xi}$ belongs to $A$ and is
the smallest projection $p\in A$ satisfying $p\xi=\xi$; this is precisely
$s_A(\omega_\xi)$.  The set
$D\cap(1-e)K$ is $A'$-invariant.  Hence
\[
 W_e=\overline{\operatorname{span}}
       \{\overline{A'\xi}:\xi\in D,\ e\xi=0\},
\]
whose projection is the join of the corresponding support projections.
Substitution in \eqref{eq:re} proves \eqref{eq:main}.  Finally,
$e=e^{\mathrm{hsc}}$ is equivalent to $r_e=1-e$, which is exactly
\eqref{eq:fixed}.
\end{proof}

\section{The full-factor and finite-corank cases}

\begin{corollary}\label{cor:factor}
Suppose $M=\mathcal B(H)$ in its defining representation.  Then
$A=\mathcal B(K)$ and, for every projection $e$ on $K$,
\begin{equation}\label{eq:factor}
 e^{\mathrm{hsc}}
 =1-P_{\overline{\ker(e)\cap\TC(H)}^{\,\|\cdot\|_2}}.
\end{equation}
If $F=\ker e$ is finite-dimensional, then
\begin{equation}\label{eq:finite}
                  e^{\mathrm{hsc}}=1-P_{F\cap\TC(H)}.
\end{equation}
In particular, when $F=\mathbb C\xi$,
\[
 e^{\mathrm{hsc}}=
 \begin{cases}
 e,&\xi\in\TC(H),\\
 1,&\xi\notin\TC(H).
 \end{cases}
\]
\end{corollary}

\begin{proof}
For $M=\mathcal B(H)$, the standard action gives
$M\bar\otimes M^{\mathrm{op}}=\mathcal B(\HS(H))$; thus
Theorem~\ref{thm:formula} immediately gives \eqref{eq:factor}.  If $F$ is
finite-dimensional, its linear subspace $F\cap\TC(H)$ is closed in $F$,
so the closure in \eqref{eq:factor} may be removed.  In dimension one the
intersection is either all of $F$ or $\{0\}$.
\end{proof}

The source's Example~4.12 is the second alternative: its normal vector is
the diagonal Hilbert--Schmidt operator with diagonal $(1/n)$, which is not
trace class, and the closure fills the entire diagonal defect.  The
corollary shows that this phenomenon is universal for every corank-one
projection, not special to a diagonal model.

\section{Commutative interpretation and remaining gap}

For the standard representation $M=L^\infty(X)\subseteq\mathcal B(L^2X)$,
the algebra $A$ identifies with $L^\infty(X\times X)$ on
$K=L^2(X\times X)$.  If $e=1_E$, formula~\eqref{eq:main} says that
$1-e^{\mathrm{hsc}}$ is the indicator of the essential union of supports
of all trace-class kernels supported in $E^c$.  Here the trace-class
kernels are the image of
\(L^2(X)\widehat\otimes_\pi\overline{L^2(X)}\) in $L^2(X\times X)$.
For a countable atomic $X$, the matrix units show that every projection is
fixed.  In the diffuse case the fixed-point problem is a form of operator
synthesis: which measurable complements are generated, in $L^2$, by the
trace-class kernels they support?

Thus Theorem~\ref{thm:formula} is an exact projection-lattice description
and gives nontrivial complete special cases.  What remains open is a
representation-free characterization of the join in \eqref{eq:main} using
only intrinsic structure of an arbitrary $M\bar\otimes M^{\mathrm{op}}$,
or an equally explicit classification of its fixed projections.  The
paper's separate abstract Hilbert-module reconstruction question for
operator-valued weights is not addressed here.

\section{Novelty note}

The proof is a short consequence of Proposition~4.11 once support
projections are brought into the picture.  The source does not state
formula~\eqref{eq:main} or the finite-corank classification.  Searches of
the local run corpus through 11 August 2026 found no duplicate result, but
the argument may be regarded as an implicit reformulation by specialists.
Novelty confidence is therefore moderate, while mathematical-validity
confidence is high.  Independent expert review is requested.

\begin{thebibliography}{9}
\bibitem{Daws}
M.~Daws,
\emph{Quantum graphs in infinite-dimensions: Hilbert--Schmidts and Hilbert
modules}, arXiv:2511.23121v1 (2025),
\url{https://arxiv.org/abs/2511.23121}.
\end{thebibliography}

\end{document}
